\documentclass[10pt,a4paper]{amsart}
\usepackage[margin=19mm]{geometry}
\usepackage{amsmath,amssymb,mathtools}
\usepackage[scr=rsfs,cal=euler]{mathalfa}
\usepackage{xfrac}
\usepackage[unicode,hidelinks,bookmarks=false]{hyperref}
\newcommand{\ie}{i.e.\ }
\newcommand{\cf}{cf.\ }
\newcommand{\resp}{respectively\ }
\newcommand{\Subsection}{\subsection}
\providecommand{\nobreakdash}{\nobreak\hbox{-}}
\setlength{\emergencystretch}{2em}
\multlinegap0pt
\usepackage{amssymb}
\usepackage{mathtools}
\newtheorem{theorem}{Theorem}
\newtheorem{lemma}{Lemma}
\newcommand{\Hin}{\mathrm{Hin}}
\newcommand{\p}{\mathbb{P}}
\title{Random interlacements on transient weighted graphs: 0-1 laws and FKG inequality}
\author{Orph\'ee Collin}
\date{Source: arXiv:2603.10896v1 (CC BY 4.0)}
\begin{document}
\maketitle

\noindent\emph{Source and licence.} Random interlacements on transient weighted graphs: 0-1 laws and FKG inequality. Version arXiv:2603.10896v1 (CC BY 4.0), distributed by arXiv under the Creative Commons Attribution 4.0 International licence, http://creativecommons.org/licenses/by/4.0/, as recorded in the arXiv licence metadata for this exact version and shown on the abstract page https://arxiv.org/abs/2603.10896v1. Full licence notice: \texttt{licenses/arxiv-2603.10896-CC-BY-4.0.txt}.\par\smallskip

\noindent\textit{Editorial scope.} Counted unit: the article's theorem th:quantitative\_criterion (source0.tex lines 686--691). The article's complete original proof of it is delivered with it (source0.tex lines 1359--1446, one \texttt{proof} environment of 88 lines, which the article places away from the statement). No mathematics is written, repaired or re-derived: the statement and the proof are the article's own lines, in the article's own notation, and no retained line is substituted or re-flowed.  Retained uncounted as the source-local context the proof needs: lines 1244--1251; lines 1275--1286.  Nothing was omitted from any retained span, so the package carries no omission record.  No retained line contains a citation, so the wrapper carries no bibliography and no bibliography entry.  No retained line contains a figure environment, an image-inclusion command or drawing code, so no image is delivered and the package declares no asset dependency.  Editorial formatting only, in addition: an AMS \texttt{amsart} preamble that restates the article's own declarations and macros used by the retained lines, namely the environments \texttt{theorem}, \texttt{lemma} on separate counters, together with the standard packages those lines need.  The retained environments print in this document as Theorem 1, Lemma 1 in order of appearance, whereas the article prints the same items with the numbering of its own full document.  The complete native source of the article as distributed in the arXiv e-print for this exact version is bundled unchanged as \texttt{sources/196174/source0.tex}.
\begin{lemma}\label{lem:criterion_dTV_to_0_law}
  The 0-1 law for $\mathcal{T}_{RI}$ holds if and only if for every finite vertex set $K$, for every admissible hinge configuration $\rho$ on $K$, as $L\to \mathbb{V}$,
  \begin{equation}\label{eq:criterion_d_TV_cond}
    d_{TV}\left(\mathbb{P}(\Hin_L \in \cdot | \Hin_K=\rho),
    \mathbb{P}(\Hin_L \in \cdot | \Hin_K=0) \right) \to 0.
  \end{equation}
  Furthermore, the condition can be restricted to $\rho$'s that consist of a single hinge couple.
\end{lemma}

\begin{lemma}\label{lem:general_cond_dTVto0_for_PPP}
For every $n\ge 1$ we consider a finite set $\mathcal{X}_n$, and two finite measures $\nu_n$ and $\pi_n$ on $\mathcal{X}_n$, $\pi_n$ being a probability measure. Then, 
\begin{equation}
  d_{TV}(PPP(\nu_n) \oplus \pi_n  , PPP(\nu_n) ) \underset{n\to\infty}{\longrightarrow} 0
\end{equation}
if and only if, for every $\epsilon>0$, 
\begin{equation}
  \sum_{x\in \mathcal{X}_n} 
  \left(\pi_n(x)-\epsilon \nu_n(x)\right)_+  
  \underset{n\to\infty}{\longrightarrow} 0\, .
\end{equation}
\end{lemma}

\begin{theorem}\label{th:quantitative_criterion}
  The 0-1 law for $\mathcal{T}_{RI}$ holds if and only if for every site $x\in \mathbb{V}$, for every $\epsilon>0$, as $L\to \mathbb{V}$,
  \begin{equation}
    \sum_{x', y'\in \partial_X L} h_L(x', y') \left(P_{x'}[\tau_x < \infty | X_{\lambda_L}=y']-\epsilon \right)_+ \to 0.
  \end{equation}
\end{theorem}

\begin{proof}[Proof of Theorem \ref{th:quantitative_criterion}]
Assume that the 0-1 law for $\mathcal{T}_{RI}$ holds under $\mathbb{P}$. Consider a vertex $x$ in $\mathbb{V}$, set $K=\{x\}$ and consider the hinge couple $(x,x)$ on $K$. Applying Lemma \ref{lem:criterion_dTV_to_0_law}, we get the convergence towards zero, as $L\to\mathbb{V}$, of the total variation distance between 
\begin{itemize}
  \item the law 
  $\mathbb{P}(\Hin_L \in \cdot | \Hin_K=0)$, that is, a Poisson point process with intensity measure 
  $h_{L}^{x}(x',y') := h_L(x',y') P_{x'}(\tau_x=\infty|X_{\lambda_L}=y')$,
  \item and the law 
  $\mathbb{P}(\Hin_L \in \cdot | \Hin_K=\delta_{(x,x)})$, that is, the law of the superposition of the above Poisson point process and an independent point sampled under the probability 
  $p_{L}^{x}(x', y'):= P_x(X_{\lambda_L}= x'| \tau_x^+=\infty) P_x(X_{\lambda_L}= y'| \tau_x^+=\infty)$.
\end{itemize}
Applying Lemma \ref{lem:general_cond_dTVto0_for_PPP} for any exhaustion $(L_n)_{n\ge 0}$ of $\mathbb{V}$, this implies that for every $\epsilon>0$, 
\begin{equation}
  \sum_{x',y' \in \partial_X L} 
  \left(p_{L}^{ x}(x', y') 
  - \epsilon h_{L}^{x}(x',y') \right)_+  \underset{L \to \mathbb{V}}\longrightarrow 0.
\end{equation}
Bounding $h_{L}^{x}$ by $h_L$, this implies that  
\begin{equation}
  \sum_{x',y'\in  \partial_X L} 
  \left(p_{L}^{x}(x', y')  - \epsilon h_L(x', y')   \right)_+  
  \underset{L \to \mathbb{V}}{\longrightarrow} 0.
\end{equation}
Now, by reversibility:
\begin{equation}
  a_x  P_x[X_{\lambda_L}=x', \tau_x^+=\infty]    P_x[X_{\lambda_L}=y']  = a_{x'} P_{x'}[\tau_L=\infty]  P_{x'}[\tau_x < \infty, X_{  \lambda_L}=y'].
\end{equation}
We also note that $P_x[X_{\lambda_L} = z,\tau_x^+=\infty]= P_x[\tau_x^+=\infty] P_x[X_{\lambda_L} = z]$, for every $z\in \mathbb{V}$.
It follows that:
\begin{equation}\label{eq:rewrite_proba_induced_hinges_singleton}
 p_{L}^x(x', y') =  \frac{h_L(x', y')}{a_x P_x(\tau_x^+=\infty) } 
  P_{x'}[\tau_x < \infty | X_{\lambda_L}=y'].
\end{equation}
Recalling that $x$ is fixed, and playing on $\epsilon$, we deduce that for every $\epsilon>0$
\begin{equation}
  \sum_{x',y' \in \partial_X L} 
  h_L(x', y')
  \left(P_{x'}[\tau_x < \infty | X_{\lambda_L}=y']
  - \epsilon  \right)_+  \underset{L \to \mathbb{V}}{\longrightarrow} 0.
\end{equation}
This is the condition of Theorem \ref{th:quantitative_criterion}.

\medskip 

Conversely, let us assume that the condition of Theorem \ref{th:quantitative_criterion} holds, and let us pick a finite vertex set $K$, and $(x,y)$ an admissible hinge couple on $K$. Following Lemma \ref{lem:general_cond_dTVto0_for_PPP}, our aim is to show the convergence towards zero, as $L\to\mathbb{V}$, of the total variation distance between
\begin{itemize}
  \item the law 
  $\mathbb{P}(\Hin_L \in \cdot | \Hin_K=0)$, that is, a Poisson point process with intensity measure 
  $h_{L}^{K}(x',y') := h_L(x',y') P_{x'}(\tau_K=\infty|X_{\lambda_L}=y')$
  \item and the law 
  $\mathbb{P}(\Hin_L \in \cdot | \Hin_K=\delta_{(x,y)})$, that is, the law of the superposition of the above Poisson point process and an independent point sampled under the probability 
  $p_{L}^{K, x,y} (x', y'):= P_x(X_{\lambda_L}= x'| \tau_K^+=\infty) P_y(X_{\lambda_L}= y'| \tau_K^+=\infty)$.
\end{itemize}
By Lemma \ref{lem:general_cond_dTVto0_for_PPP}, it is enough to establish that for every $\epsilon>0$, 
\begin{equation}\label{eq:sum_pLKxy-epshLK_to0}
  \sum_{x',y' \in \partial_X L} 
  \left(p_L^{K, x, y}(x',y') 
  - \epsilon h_L^{K}(x',y') \right)_+  \underset{L \to \mathbb{V}}{\longrightarrow} 0.
\end{equation}
By reversibility:
\begin{equation}
  \begin{aligned}
  a_x &  \p_x[X_{\lambda_L}=x', \tau_K^+=\infty]   \p_x[X_{\lambda_K}=y] \p_y[X_{\lambda_L}=y' | \tau_K^+=\infty] \\
   =  & a_{x'} \p_{x'}[\tau_L^+=\infty]
   \p_{x'}[\tau_K < \infty, X_{\tau_K} = x, X_{\lambda_K}=y, X_{\lambda_L}=y']
  \end{aligned}
\end{equation}
hence
\begin{equation}
  p_L^{K, x, y}(x',y')= \frac {h_L(x', y')}{ h_K(x, y)} 
  \p_{x'}[\tau_K < \infty, X_{\tau_K} = x, X_{\lambda_K}=y| X_{\lambda_L}=y'],
\end{equation}%\label{eq:rewrite_proba_induced_hinges}
and therefore, playing on $\epsilon$, \eqref{eq:sum_pLKxy-epshLK_to0} is equivalent to the fact that for every $\epsilon>0$, 
\begin{equation}
  \begin{aligned}
    \sum_{x', y'\in \partial_X L} &  h_L(x', y') \\
    &\times \big(  P_{x'}[\tau_K<\infty, X_{\tau_K}=x, X_{\lambda_K}=y |X_{\lambda_L}=y']  -\epsilon P_{x'}(\tau_K=\infty|X_{\lambda_L}=y') \big)_+ 
  \end{aligned}
  \end{equation}
vanishes as $L \to \mathbb{V}$. Bounding $P_{x'}[\tau_K<\infty, X_{\tau_K}=x, X_{\lambda_K}=y  |X_{\lambda_L}=y']$ by $P_{x'}[\tau_K<\infty|X_{\lambda_L}=y']$, the latter follows from the fact that for every $\epsilon>0$,
\begin{equation}
    \sum_{x',y' \in \partial_X L} h_L(x', y') \left((1+\epsilon) P_{x'}[\tau_K<\infty | X_{\lambda_L}=y']-\epsilon \right)_+ \underset{L \to \mathbb{V}}{\longrightarrow} 0,
  \end{equation}
which, finally, can be derived elementary from the condition of Theorem \ref{th:quantitative_criterion} by using the bound:
\begin{equation}
  P_{x'}[\tau_K<\infty|X_{\lambda_L}=y'] \le \sum_{z\in K}  P_{x'}[\tau_z<\infty|X_{\lambda_L}=y']
\end{equation}
and playing on $\epsilon$.
\end{proof}

\end{document}
